\color{unidarkgreen}{\subsection{Максимальная токовая защиты стороны НН1(НН)}}
\color{black}


\begin{enumerate}

\item
Максимальная токовая защита (МТЗ) стороны НН1(НН) предназначена для резервирования действия основных быстродействующих защит. Защита выполнена двухступенчатой. Измерительные цепи защиты включены на токи фаз А, В, С со стороны НН1(НН) трансформатора. 

\item
Каждая из ступеней имеет регулируемые значения по току и времени срабатывания. Значения срабатывания по току определяются параметрами \textbf{«Iср МТЗ-1 НН1(НН)»} и \textbf{«Iср МТЗ-2 НН1(НН)»} для
первой и второй ступеней соответственно. Задержки на срабатывание защиты задаются параметрами \textbf{«Тср МТЗ-1 НН1(НН)»} и \textbf{«Тср МТЗ-2 НН1(НН)»}.

\item
Каждая из ступеней защиты может выполнена с пуском по напряжению. Ввод контроля пуска по напряжению в логику работы ступеней осуществляется программными накладками \textbf{«Режим КПОН МТЗ-1 НН1(НН)»} и \textbf{«Режим КПОН МТЗ-2 НН1(НН)»} для первой и второй ступеней соответственно.

\item
Каждая из ступеней защиты может выполнена с блокировкой по второй гармонике (отношению уровня второй гармоники к первой). Данная блокировка реализована на базе модуля «ОВ БНТ НН1(НН)». Ввод контроля броска намагничивающего тока программными накладками \textbf{«Режим БНТ МТЗ-1 НН1(НН)»}, \textbf{«Режим БНТ МТЗ-2 НН1(НН)»}. 

\item
Для обеспечения несрабатывания МТЗ НН1(НН) при броске тока намагничивания используется перекрестная блокировка по 2 гармонике в фазных токах НН1(НН). Вводится программной накладкой \textbf{«БНТ МТЗ НН1(НН)»}.  Блокировка основана на контроле отношения действующего значения второй гармонике к первой токов НН1(НН). Данная блокировка вводиться в работу для каждой ступени независимо с помощью уставок \textbf{«Режим БНТ МТЗ-1 НН1(НН)»} и \textbf{«Режим БНТ МТЗ-2 НН1(НН)»}. Имеется возможность регулировать порог срабатывания блокировки, с помощью уставки \textbf{«I2г/I1г МТЗ НН1(НН)»}. Блокировка МТЗ НН1(НН) по 2 гармонике выводится из действия при величине соответствующего фазного тока ниже 0,04Iном. Также блокировка МТЗ НН1(НН) по 2 гармонике выводится из действия при достижении величины действующего значения тока выше уставки \textbf{«Iср БНТ МТЗ НН1(НН)»}.

\item
Алгоритм работы МТЗ НН1(НН) приведен на рисунке \ref{pic:MTZ_NN1}. Уставки МТЗ НН1(НН) приведены в таблице \ref{app_set:mtznn1}.

\medskip
\begin{figure}[H]
\centering
\includegraphics[width=1\textwidth,height=1\textheight,keepaspectratio]{\fbpath/06.01 ДЗТ/МТЗ НН1/img/MTZ_NN1.pdf}
\captionof{figure}{Функциональная схема логики МТЗ НН1(НН)}\label{pic:MTZ_NN1}
\end{figure}

\end{enumerate}

